\documentclass[11pt]{article}

\usepackage{amsmath,amssymb}
\usepackage[margin=1in]{geometry}
\usepackage{booktabs}
\usepackage{array}
\usepackage{hyperref}

\title{Localized-Mode Protection and Failure-Mechanism Feasibility on a Frozen Branch-Local Cochain-Laplacian Hierarchy}
\author{Tomislav Rupi\'c \\ Independent Researcher}
\date{March 2026}

\begin{document}

\maketitle

\begin{abstract}
Phase XI of the frozen HAOS-IIP program asks whether the already frozen localized candidate
\texttt{low\_mode\_localized\_wavepacket} is protected by reproducible structural mechanisms or is only an accidental configuration. The phase is not allowed to alter the Phase V readout layer, the Phase VI branch-local operator hierarchy $\Delta_h$, the Phase VII spectral-feasibility manifest, the Phase VIII short-time trace contract, the Phase IX coefficient-stabilization contract, the Phase X cautious bridge bundle, or the Phase X proto-particle candidate definition. It therefore tests protection only through deterministic perturbation channels on the frozen branch and on a matched altered-connectivity control. Four probe classes are used: local stiffness defects, connectivity micro-surgery, bounded phase randomization, and boundary leakage. The branch exhibits a bounded survival basin with refinement-ordered first-failure amplitudes for local stiffness defects, increasing from $0.5$ at $n_{\mathrm{side}}=48$ to $0.7$ at $60$ and $0.9$ at $72$, with no failure observed up to amplitude $1.0$ at $84$. Phase-randomization failure is sharp and refinement-stable, with first failure at amplitude $0.08$ on every tested level. The branch persistence times grow across refinement as $\tau(h) = [1.2, 2.0, 2.4, 3.6]$, with fitted power-law exponent $q = 1.8684$ and $R^2 = 0.9774$, while the control times remain shorter, $[0.4,1.2,1.6,2.0]$. Failure channels are consistent: strong local stiffness defects are spectral-mixing dominated, connectivity micro-surgery is topological-defect triggered, phase randomization is incoherence triggered, and the scanned boundary leakage channel remains inside the persistent basin. The result is bounded: Phase XI establishes localized-mode protection and failure-mechanism feasibility for the frozen operator hierarchy. It does not assert particles, fields, or new physical interpretation.
\end{abstract}

\section{Scope and Frozen Contracts}

Phase XI is a protection and failure-mode diagnostic only. It does not introduce a new operator, a new normalization, or a new interpretive layer. It consumes the following frozen authority contracts without change:
\begin{itemize}
\item the Phase V recovery-statistics readout layer;
\item the Phase VI branch-local periodic DK2D block cochain Laplacian hierarchy $\Delta_h$;
\item the Phase VII spectral-feasibility outputs;
\item the Phase VIII short-time trace window and probe grid;
\item the Phase IX coefficient stabilization diagnostics;
\item the Phase X cautious continuum-bridge feasibility bundle;
\item the Phase X proto-particle candidate definition.
\end{itemize}

The concrete artifact references used here are:
\begin{itemize}
\item \url{phase6-operator/phase6_operator_manifest.json};
\item \url{phase7-spectral/phase7_spectral_manifest.json};
\item \url{phase8-trace/phase8_trace_manifest.json};
\item \url{phase9-invariants/phase9_manifest.json};
\item \url{phase10-bridge/phase10_manifest.json};
\item \url{phaseX-proto-particle/phaseX_integrated_manifest.json};
\item \url{phase11-protection/phase11_manifest.json}.
\end{itemize}

Nothing in this phase modifies earlier manifests or \texttt{haos\_core}. All logic is isolated inside \texttt{phase11-protection/}.

\section{Protection Program}

The frozen candidate is the Phase X refinement-stable branch excitation
\[
\texttt{low\_mode\_localized\_wavepacket},
\]
evaluated on the refinement-ordered branch hierarchy
\[
h = \frac{1}{n_{\mathrm{side}}}, \qquad
n_{\mathrm{side}} \in \{48,60,72,84\}.
\]

The phase uses one fixed evaluation time
\[
\tau = 0.8,
\]
and one fixed persistence-time grid
\[
\tau \in \{0.0, 0.4, 0.8, 1.2, 1.6, 2.0, 2.4, 3.0, 3.6, 4.2\}.
\]

Four deterministic perturbation classes are scanned:
\begin{enumerate}
\item local stiffness defect in a compact neighborhood;
\item connectivity micro-surgery with fixed edge rewiring/removal pattern;
\item bounded phase-randomization window on the initial state;
\item boundary leakage probe implemented as deterministic dissipative coupling outside a radius threshold.
\end{enumerate}

For each perturbation, Phase XI records localization width, concentration, participation growth, overlap with the unperturbed branch excitation, a recovery-style score, a persistence class, and a failure-mechanism label.

\section{Bounded Survival Basin}

The clearest refinement-ordered protection result comes from the local stiffness defect probe.

\begin{table}[t]
\centering
\small
\begin{tabular}{@{}rrrr@{}}
\toprule
$n_{\mathrm{side}}$ & $h$ & max persistent amplitude & first failure amplitude \\
\midrule
48 & 0.020833 & 0.4 & 0.5 \\
60 & 0.016667 & 0.6 & 0.7 \\
72 & 0.013889 & 0.8 & 0.9 \\
84 & 0.011905 & 1.0 & none observed up to 1.0 \\
\bottomrule
\end{tabular}
\caption{Branch local-stiffness survival basin.}
\label{tab:stiffness_basin}
\end{table}

This gives a refinement-ordered first-failure sequence
\[
0.5,\;0.7,\;0.9,\;>1.0,
\]
which is the main protection-basin result of the phase.

The phase-randomization window shows a different kind of structure. Its threshold is not widening with refinement; instead it is sharply locked:
\[
\text{first failure amplitude} = 0.08
\]
on every tested level. This indicates a narrow coherence-sensitive channel rather than a broad structural damage channel.

The branch mean survival fractions across the full scanned grids are:
\begin{itemize}
\item local stiffness defect: $0.7273$;
\item connectivity micro-surgery: $0.9583$;
\item phase randomization: $0.4444$;
\item boundary leakage: $1.0$.
\end{itemize}

The boundary leakage result is deliberately bounded: it does not claim immunity. It states only that over the declared amplitude range the scanned leakage channel did not exit the persistence basin.

\section{Failure Mechanisms}

Phase XI does not just label survival or failure. It also tags the dominant failure channel for each perturbation family. The branch produces a consistent map:
\begin{itemize}
\item local stiffness defect: spectral-mixing dominated;
\item connectivity micro-surgery: topological-defect triggered;
\item phase randomization window: incoherence triggered;
\item boundary leakage probe: survives on the scanned range.
\end{itemize}

This is the main failure-mechanism result. Breakdown does not appear random. Different perturbation classes fail through distinct and reproducible channels.

\section{Persistence-Time Scaling}

Phase XI defines the candidate persistence time as the largest value on the declared $\tau$ grid for which the excitation remains inside the frozen scan-level compactness gate. On the branch, the persistence times are
\[
\tau(h) = [1.2,\;2.0,\;2.4,\;3.6].
\]
On the deterministic control they are
\[
\tau_{\mathrm{ctrl}}(h) = [0.4,\;1.2,\;1.6,\;2.0].
\]

\begin{table}[t]
\centering
\small
\begin{tabular}{@{}rrrr@{}}
\toprule
$n_{\mathrm{side}}$ & branch $\tau$ & control $\tau$ & branch/control ratio \\
\midrule
48 & 1.2 & 0.4 & 3.0 \\
60 & 2.0 & 1.2 & 1.666667 \\
72 & 2.4 & 1.6 & 1.5 \\
84 & 3.6 & 2.0 & 1.8 \\
\bottomrule
\end{tabular}
\caption{Persistence-time comparison between branch and control.}
\label{tab:tau_scaling}
\end{table}

The branch fit to a minimal power law
\[
\tau(h) \sim h^{-q}
\]
gives
\[
q = 1.868386844019,
\qquad
R^2 = 0.977397586681.
\]
The control fit is less clean and stays below the branch at every tested level:
\[
q_{\mathrm{ctrl}} = 2.819218651794,
\qquad
R^2 = 0.907048419900.
\]

The Phase XI success criterion does not require a universal exponent. It requires structured scaling and clear branch/control separation. Those conditions are met.

\section{Structural Correlates}

The phase also tests whether simple structural diagnostics track survival and persistence. The strongest branch correlate is the coefficient-ratio shift proxy inherited from the frozen spectral-trace stack:
\[
\mathrm{corr}(\tau,\text{ratio-shift proxy}) = -0.930992601909.
\]

The operator-survival indicator also shows a nontrivial relation to local structure:
\begin{itemize}
\item branch survival vs local spectral density proxy: $0.492373478603$;
\item branch survival vs local connectivity variance: $-0.804337994826$.
\end{itemize}

These are not interpretive invariants. They are descriptive correlates showing that protection is not numerically random with respect to the already frozen branch descriptors.

\section{Control Comparison}

The deterministic altered-connectivity control is tested under the same perturbation grids and the same classification rules. Its mean survival fractions are lower:
\begin{itemize}
\item local stiffness defect: $0.5$;
\item connectivity micro-surgery: $0.5$;
\item phase randomization: $0.2222$;
\item boundary leakage: $0.5$.
\end{itemize}

At the coarser refinement levels $n_{\mathrm{side}} = 48$ and $60$, the control leaves the persistent basin immediately under several perturbation families, with first failure amplitude $0.0$. This is the key qualitative separation. The branch has a measurable protection basin where the control does not.

\section{Bounded Interpretation}

The Phase XI result is narrow by design. It supports the following claim only: within the already frozen operator and spectral-trace contracts, the localized branch candidate has a deterministic survival basin, its breakdown channels are class-specific and reproducible, and its persistence time grows across refinement more strongly than the matched control.

This is not a particle claim. It is not a field-dynamics claim. It is not a continuum statement. It is not a new operator program. It is only a protection and failure-mechanism feasibility result on a frozen hierarchy.

\section{Conclusion}

Phase XI identifies a bounded protection basin for the frozen localized mode, with refinement-ordered local-stiffness failure thresholds, a sharp and repeatable incoherence threshold under phase randomization, structured persistence-time growth, and a clear deterministic-control separation. The artifact bundle is frozen in:
\begin{itemize}
\item \url{phase11-protection/phase11_manifest.json};
\item \url{phase11-protection/phase11_summary.md};
\item \url{phase11-protection/runs/phase11_perturbation_survival_ledger.csv};
\item \url{phase11-protection/runs/phase11_failure_classification.csv};
\item \url{phase11-protection/runs/phase11_persistence_scaling.csv};
\item \url{phase11-protection/runs/phase11_runs.json}.
\end{itemize}

Phase XI establishes localized-mode protection and failure-mechanism feasibility for the frozen operator hierarchy.

\end{document}
